\documentclass[12pt,a4paper]{article}
\usepackage{BerichtVorlage-erweitert}		% Stil einbinden

\title{Interaktive Lernplattform für Grundlagen Informatik}
\author{M. Nöth, R. Schröder, E. Ebert, D. T. Pham, K. Popov}
\date{\today}

% Optional: Header / Footer anpassen
% Standardmäßig steht der Autor mittig im Header
% und die Seitenzahl mittig im Footer.
% \setheaderleft{abc}
% \setheadercenter{Autor}
\setheaderright{Projektarbeit}
% \setfooterleft{Entwurf}
% \setfootercenter{Seite}
% \setfooterright{Vertraulich}

\setheaderleft{\hspace*{-0.68cm}\includegraphics[width=2cm]{THWS.pdf}}

\addbibresource{Literatur.bib}

% Beispiel-Akronyme für das Abkürzungsverzeichnis.
\newacronym[
	description={Hochschule aus Würzburg und Schweinfurt (Bayern) \cite{thws2026}}	% Geht nur wegen doppeltem biber Lauf in compile-erweitert.ps1 -> sonst vorher per \nocite{thws2026} laden
]{thws}{THWS}{Technische Hochschule Würzburg-Schweinfurt}

\newacronym[
	description={Offener Standard, über den Sprachmodelle externe Werkzeuge und Datenquellen anbinden}
]{mcp}{MCP}{Model Context Protocol}

\newacronym[
	description={Protokoll zwischen Editor und Sprachserver für Codenavigation, Typinformationen und Refactoring}
]{lsp}{LSP}{Language Server Protocol}

\glsunsetall    % only show short forms of all acronyms

\begin{document}

\begin{Cover}
	\begin{center}
		\includegraphics[width=10cm]{THWS.pdf}

		\vspace*{5em}

		{\Large \textbf{\thetitle}}

		\vspace*{3em}

		Projektarbeit -- SS 26

		\vspace*{2em}

		\thedate

		\vspace*{25em}

		\noindent\begin{tabular}{@{}lllr}
			\textbf{Studenten}:	& Markus Nöth		&\href{mailto:markus.noeth@study.thws.de}{markus.noeth@study.thws.de}			& 5123027\\[2ex]
								& Richard Schröder	&\href{mailto:richard.schroeder@study.thws.de}{richard.schroeder@study.thws.de}	& 5123042\\[2ex]
								& Enrico Ebert		&\href{mailto:enrico.ebert@study.thws.de}{enrico.ebert@study.thws.de}			& 5123098\\[2ex]
								& Duy Tan Pham		&\href{mailto:duytan.pham@study.thws.de}{duytan.pham@study.thws.de}				& 5123063\\[2ex]
								& Kristian Popov	&\href{mailto:kristian.popov@study.thws.de}{kristian.popov@study.thws.de}		& 5123029\\[4ex]
			\textbf{Dozent}:	& Prof. Dr. Peter Braun
		\end{tabular}
	\end{center}
\end{Cover}

\begin{PreContent}
	\tableofcontents
	\listoffigures
	\listoftables
\end{PreContent}

\begin{Content}
	\subNsection{0}{Einleitung}
Die vorliegende Projektarbeit entstand im Sommersemester 2026 an der \gls{thws} unter Betreuung von Prof.~Dr.~Braun. Ziel des Projekts war die Konzeption und Umsetzung einer interaktiven Lernplattform f{\"u}r das Modul \glqq Grundlagen der Informatik\grqq{}.

\subNsection{1}{Problemstellung}{Author}

Die Hürde der Abstraktion in Grundlagen Informatik

\subNsection{1}{Zielsetzung der Arbeit}{Author}

Didaktische Aufbereitung komplexer Modelle durch interaktive Simulationen

\subNsection{1}{Methodik}{Markus Nöth}

Vorgehensweise im Team und Rollenverteilung


	\subNsection{0}{Übergreifende Architektur \& Systemdesign}
\subNsection{1}{Technologiestack}{Author}

React 19, TypeScript, Vite, Tailwind CSS 4, shadcn/ui, Zustand, i18next, Lucide Icons

\subNsection{1}{Modulare Softwarearchitektur}{Author}

Single Source of Truth (\verb|moduleCatalog.tsx|), lazy Loading, Shared Topic Pattern – Sicherstellung der Erweiterbarkeit für neue Themen

\subNsection{1}{Unified UI-Konzept}{Duy Tan Pham}

Wiederverwendbare Shared Components, konsistentes Designsystem (CSS-Variablen, Dark Mode, moduleigene Akzentfarben, Geist Font)

\subNsection{1}{Simulator-Architektur}{Richard Schröder}

\begin{figure}[H]
	\centering
	\begin{tikzpicture}[
		font=\footnotesize,
		box/.style={
			draw,
			rounded corners=2pt,
			align=center,
			inner sep=4.5pt
		},
		umlclass/.style={
			draw,
			rounded corners=2pt,
			align=left,
			inner sep=5pt,
			fill=black!8
		},
		flow/.style={
			-{Stealth[length=2.5mm,width=1.8mm]},
			semithick
		},
		realize/.style={
			-{Triangle[open,length=2.8mm,width=2.5mm]},
			dashed,
			semithick
		},
		lbl/.style={
			font=\scriptsize\itshape,
			inner sep=2pt
		}
	]
		% 1. Ebene: UI-Bedienelemente
		\node[box, fill=black!4, text width=7.2cm] (controls) at (0, 0) {
			\textbf{\guillemotleft Component\guillemotright}\\[1pt]
			\textbf{SimulatorControls}\\[2pt]
			\scriptsize Play, Pause, Schritt vor/zurück, Reset, Tempo
		};

		% 2. Ebene: React-Hook / Steuerung
		\node[box, fill=black!4, text width=7.2cm, below=0.65cm of controls] (hook) {
			\textbf{\guillemotleft Hook\guillemotright}\\[1pt]
			\textbf{\texttt{useSimulationEngine}}\\[2pt]
			\scriptsize Timer (\texttt{setInterval}), Status (\texttt{RUNNING}, \dots), React-State-Sync
		};

		% 3. Ebene: Zentrales Interface (Fokus auf Kernmethoden)
		\node[umlclass, text width=7.2cm, below=0.65cm of hook] (isim) {
			\centering\textbf{\guillemotleft interface\guillemotright}\\[1pt]
			\centering\textbf{\texttt{ISimulator<TConfig, TState>}}\\[2pt]
			\rule{\linewidth}{0.4pt}\\[2pt]
			\scriptsize
			\texttt{+ initialize(config, state?): void}\\[1pt]
			\texttt{+ step() / stepBack() / reset(): void}\\[1pt]
			\texttt{+ getCurrentState(): TState}\\[1pt]
			\texttt{+ isFinished(): boolean}
		};

		% 4. Ebene: Konkrete Fachsimulatoren (exakt symmetrisch und entzerrt)
		\node[box, fill=black!4, text width=3.4cm, below=1.1cm of isim.south, xshift=-2.05cm] (fa) {
			\textbf{\texttt{FiniteAutomaton-}}\\\textbf{\texttt{Simulator}}\\[2pt]
			\scriptsize DEA / NEA, Wortindex
		};
		\node[box, fill=black!4, text width=3.4cm, left=0.4cm of fa] (tm) {
			\textbf{\texttt{TuringMachine-}}\\\textbf{\texttt{Simulator}}\\[2pt]
			\scriptsize Band, Kopf, $\delta$-Tabelle
		};
		\node[box, fill=black!4, text width=3.4cm, below=1.1cm of isim.south, xshift=2.05cm] (moore) {
			\textbf{\texttt{MooreMachine-}}\\\textbf{\texttt{Simulator}}\\[2pt]
			\scriptsize Zustand, Ausgabe $\lambda(s)$
		};
		\node[box, fill=black!4, text width=3.4cm, right=0.4cm of moore] (mealy) {
			\textbf{\texttt{MealyMachine-}}\\\textbf{\texttt{Simulator}}\\[2pt]
			\scriptsize Zustand, Ausgabe $\lambda(s, a)$
		};

		% Verbindungen zwischen den Ebenen mit seitlichen Beschriftungen
		\draw[flow] (controls) -- (hook)
			node[lbl, midway, right=3pt] {Aktionen \& Tempo};

		\draw[flow] (hook) -- (isim)
			node[lbl, midway, right=3pt] {steuert (\texttt{<<uses>>})};

		% Realisierung: sauberer Bus im Freiraum zwischen Ebene 3 und 4
		\path (isim.south) -- (0,0 |- fa.north) coordinate[midway] (bus);
		\draw[dashed, semithick] (tm.north) |- (bus);
		\draw[dashed, semithick] (fa.north) |- (bus);
		\draw[dashed, semithick] (moore.north) |- (bus);
		\draw[dashed, semithick] (mealy.north) |- (bus);
		\draw[realize] (bus) -- (isim.south);

	\end{tikzpicture}
	\caption[Architektur der Simulationsengine]{%
		Architektur der Simulationsengine. Die UI-Komponente \texttt{SimulatorControls} interagiert mit dem Hook \texttt{useSimulationEngine}. Dieser steuert über das gemeinsame Interface \texttt{ISimulator} die Ausführung und synchronisiert den Laufzeitzustand, während die modulspezifischen Simulator-Klassen die fachlichen Zustandsübergänge implementieren.}
	\label{fig:simulator_architektur}
\end{figure}

Turingmaschinen, endliche Automaten sowie Moore- und Mealy-Automaten verwenden dieselbe Simulationsengine. Ihre Klassen implementieren \verb|ISimulator<TConfig, TState>|. Das Interface umfasst Initialisierung, Schritt, Rückschritt, Reset, Konfiguration, aktuellen Zustand und eine Prüfung ob der Haltezustand erreicht wurde. Es verbindet damit die fachliche Simulationslogik mit der Steuerung, ohne eine konkrete Darstellung vorzuschreiben.

Die Gesamtausführung folgt dabei dem in Abbildung~\ref{fig:simulator_architektur} dargestellten Schichtenmodell: Der Hook \verb|useSimulationEngine| verwaltet Ausführungsstatus, Geschwindigkeit, Browser-Timer und Fehlerbehandlung. Nach einem Schritt übernimmt er den Zustand des Simulators in den React-State. Alle vier Maschinentypen erhalten über den Hook und die gemeinsamen Simulator-Controls Play, Pause, Schritt, Rückschritt, Reset und Geschwindigkeitssteuerung, während ihre Klassen die jeweiligen Zustandsübergänge ausführen.

Ein weiterer Maschinentyp benötigt eine Implementierung des Interfaces, einen modulspezifischen Hook-Aufruf und die Anbindung an eine Laufzeitansicht. Die allgemeine Timerlogik und die vorhandenen Simulator-Controls können dabei unverändert bleiben.


\subNsection{1}{Datenmodellierung und Transfer}{Richard Schröder}

Die Simulatoren beschreiben ihre fachlichen Daten mit TypeScript-Typen und trennen diese von Laufzeit- und Darstellungszuständen:

\begin{itemize}
    \item \textbf{Automaten:} \verb|TMProgram|, \verb|FAProgram|, \verb|MooreProgram| und \verb|MealyProgram| enthalten die editierbare Konfiguration, etwa Zustände, Alphabet, Übergänge, Startzustand und Eingabewort. Die Übergänge sind als verschachtelte \verb|Record|-Strukturen modelliert; die Namen der Zustände und Symbole sind Zeichenketten.
    \item \textbf{Ausführung:} Der jeweilige Laufzeittyp, beispielsweise \verb|TMRuntime|, \verb|FARuntime|, \verb|MooreRuntime| oder \verb|MealyRuntime|, hält den aktuellen Zustand, die Position im Eingabewort beziehungsweise auf dem Band, die ausgegebenen Symbole und den Ablaufverlauf. Die fachliche Schrittausführung arbeitet auf diesen Daten; temporäre UI-Zustände wie ein geöffneter Dialog oder eine laufende Animation bleiben in den React-Komponenten.
    \item \textbf{Übungs- und Umwandlungsaufgaben:} Das DNF-Modul hält je nach Aufgabe einen Ausdruck, eine Wahrheitstabelle oder einen Bearbeitungsstand mit Antworten getrennt vor. Der Zahlensystem-Simulator verwendet \verb|NumberSystemState| mit \verb|activeMode| und zustandsspezifischen Feldern für Ganzzahlen, negative Zahlen, Brüche, IEEE~754 und Arithmetik; der Java-Simulator überträgt den Quelltext als Zeichenkette.
\end{itemize}

Das generische Interface \verb|ISimulatorTransferCodec<TValue>| definiert die Serialisierung und Deserialisierung eines Simulationswertes. \verb|exportValue| wandelt einen Wert des Typs \verb|TValue| in einen String um; \verb|importValue| rekonstruiert ihn aus einem String. Beide Methoden dürfen ihr Ergebnis unmittelbar oder asynchron als \verb|Promise| liefern. Die daraus entstehenden Zeichenketten werden für den manuellen Import und Export sowie für die Speicherung von Training-Center-Fortschritten verwendet. Ihr Format legt die jeweilige Implementierung fest; es reicht von Base64URL-kodiertem Text bis zu einem kompakten Bitstrom.

\begin{itemize}
    \item \textbf{Gemeinsame Verwendung:} \verb|SimulatorStringTransferBox| stellt ein Eingabefeld sowie Schaltflächen zum Erzeugen, Kopieren und Laden eines Transfer-Strings bereit. Fehler beim Export oder Import werden in der Komponente abgefangen und angezeigt.
    \item \textbf{Training Center:} \verb|useTrainingCenterAutoSave| exportiert den aktuellen Wert standardmäßig alle zehn Sekunden und zusätzlich beim Aufräumen des Hooks. \texttt{exportAndStoreTrainingCenter\allowbreak Value} speichert den Payload zusammen mit einem Bezeichner; \texttt{createTrainingCenter\allowbreak TransferCodec} verpackt einen fachlichen Codec so, dass ein exportierter String diesen Speicherwert enthält und beim Import der Bezeichner geprüft wird.
    \item \textbf{Binärstrom:} Mehrere Codecs verwenden \verb|BitWriter| und \verb|BitReader|. Die resultierenden Bits werden als Base64URL ausgegeben; beim Import prüfen die Implementierungen Versionsfelder, Größenlimits, IDs und verbleibende Füllbits.
\end{itemize}

Die folgenden zwei Codecs zeigen unterschiedliche Transferstrategien:

\begin{itemize}
    \item \verb|tmProgramTransferCodec| normalisiert Turingmaschinenprogramme, entfernt ungenutzte Zustände und Symbole und vergibt kompakte numerische IDs. Für Symbole wählt der Codec zwischen einer festen Registry und UTF-8; für Übergänge vergleicht er dichte und dünn besetzte Kodierungen, jeweils mit oder ohne Zustands- und Symboldictionary, und verwendet die kürzeste Variante.
    \item \verb|javaCodeTransferCodec| kodiert Java-Quelltext als Base64URL mit einem Präfix für direkte Tokens, deflatierte Tokens oder deflatiertes UTF-8. Häufige Java-Ausdrücke und erlaubte Sonderzeichen erhalten einzelne Tokens; anschließend wählt der Codec unter den verfügbaren Varianten die kürzeste Darstellung und prüft beim Import Größenlimits sowie das Token- und UTF-8-Format.
\end{itemize}


	\subNsection{0}{Einsatz von Künstlicher Intelligenz im Projekt}
\subNsection{1}{KI Reflexion}{Markus Nöth}

Nutzung von Gemini 3.1 Pro, ChatGPT und Codex als Entwicklungswerkzeuge – Effizienzsteigerung bei der Implementierung mathematischer Logik und UI-Komponenten

\subNsection{2}{KI als Co-Developer}{Markus Nöth}

Reflektieren und so...

\subNsection{2}{Automatisierte Aufgabengenerierung}{Markus Nöth}

Reflektieren und so...

\subNsection{2}{Kritische Reflexion}{Markus Nöth}

Reflektieren und so...


\subNsection{1}{KI Reflexion}{Richard Schröder}

Nutzung von Gemini 3.1 Pro, ChatGPT und Codex als Entwicklungswerkzeuge – Effizienzsteigerung bei der Implementierung mathematischer Logik und UI-Komponenten

\subNsection{2}{KI als Co-Developer}{Richard Schröder}

Reflektieren und so...

\subNsection{2}{Automatisierte Aufgabengenerierung}{Richard Schröder}

Reflektieren und so...

\subNsection{2}{Kritische Reflexion}{Richard Schröder}

Reflektieren und so...


\subNsection{1}{KI Reflexion}{Enrico Ebert}

Nutzung von Gemini 3.1 Pro, ChatGPT und Codex als Entwicklungswerkzeuge – Effizienzsteigerung bei der Implementierung mathematischer Logik und UI-Komponenten

\subNsection{2}{KI als Co-Developer}{Enrico Ebert}

Reflektieren und so...

\subNsection{2}{Automatisierte Aufgabengenerierung}{Enrico Ebert}

Reflektieren und so...

\subNsection{2}{Kritische Reflexion}{Enrico Ebert}

Reflektieren und so...


\subNsection{1}{KI Reflexion}{Duy Tan Pham}

Nutzung von Gemini 3.1 Pro, ChatGPT und Codex als Entwicklungswerkzeuge – Effizienzsteigerung bei der Implementierung mathematischer Logik und UI-Komponenten

\subNsection{2}{KI als Co-Developer}{Duy Tan Pham}

Reflektieren und so...

\subNsection{2}{Automatisierte Aufgabengenerierung}{Duy Tan Pham}

Reflektieren und so...

\subNsection{2}{Kritische Reflexion}{Duy Tan Pham}

Reflektieren und so...


\subNsection{1}{KI-Reflexion}{Kristian Popov}

Mein Schwerpunkt lag während des Projekts auf den Themen \emph{Endliche Automaten}
und \emph{Reguläre Sprachen}.
Dabei habe ich Gemini 3.1 Pro, die kostenlose Version von ChatGPT sowie Claude Opus
4.8 und später Opus 5 ausprobiert.

Die Modelle waren dabei unterschiedlich gut.
Manche erklärten theoretische Zusammenhänge besser, während andere bei der
Implementierung oder Fehlersuche hilfreicher waren.
Überwiegend arbeitete ich mit Claude, da die Vorschläge bei Programmieraufgaben
häufiger direkt brauchbar waren und ich dadurch weniger nacharbeiten musste.

Im Laufe des Projekts ist mir außerdem aufgefallen, dass die Modelle unterschiedlich
viel Aufwand in die Planung einer Antwort stecken.
Eine höhere Planungsstufe bedeutet mehr verbrauchte Tokens und längere Wartezeiten,
führt aber nicht automatisch zu besseren Ergebnissen.
Für kleine Änderungen reichte meist eine niedrige Stufe aus, während komplexere
Aufgaben mit mehreren Dateien von einer höheren Stufe profitierten.
Deshalb habe ich je nach Aufgabe entschieden, welche Einstellung sinnvoll war.

\subNsection{2}{KI als Co-Developer}{Kristian Popov}

Die KI habe ich direkt in die Entwicklungsumgebung eingebunden.
Dadurch konnte sie auf Projektdateien zugreifen, Änderungen vorschlagen und diese
direkt im Projekt überprüfen.

Besonders hilfreich war sie bei wiederkehrenden Aufgaben wie Refactorings,
Boilerplate-Code, Unit-Tests, Internationalisierung und dem Beheben von
TypeScript-Warnungen.
Bei der Automatenlogik und beim Parsen von Grammatiken konnte sie außerdem schnell
passende Testfälle ableiten.

Bei Architekturentscheidungen ließ ich mir zunächst mehrere Lösungsansätze mit ihren
Vor- und Nachteilen erklären und traf die Entscheidung anschließend selbst.
Jede vorgeschlagene Änderung habe ich überprüft und getestet.

Unterstützt wurde diese Arbeitsweise durch Erweiterungen aus dem offiziellen
Plugin-Verzeichnis \texttt{claude-\allowbreak plugins-\allowbreak official}.
Aktiv waren mehrere, darunter \texttt{typescript-lsp} für Typinformationen über das
\gls{lsp}.
Für meine Themen am wichtigsten waren zwei:

\begin{description}[leftmargin=*,style=nextline,itemsep=0.75ex]
\item[\texttt{playwright}]
ermöglicht der KI, die Anwendung im Browser auszuführen und aus Sicht des Benutzers
zu testen.
Dadurch konnten beispielsweise Oberflächen, Navigation und Benutzerinteraktionen
überprüft werden.

\item[\texttt{chrome-devtools-mcp}]
ermöglicht über das \gls{mcp} den Zugriff auf Konsolenausgaben und das Verhalten der
laufenden Anwendung.
Dadurch konnte ich bei Fehlern genauer untersuchen, was zur Laufzeit tatsächlich
passiert.
\end{description}

Zusätzlich habe ich das Projekt einmalig mit Graphify~\cite{shamsi2026graphify}
analysiert und einen Wissensgraphen erzeugt.
Dadurch konnte die KI bei späteren Aufgaben gezielt auf relevante Dateien, Klassen
und Funktionen zugreifen, ohne das gesamte Projekt erneut einlesen zu müssen.

\subNsection{2}{Automatisierte Aufgabengenerierung}{Kristian Popov}

Für die Übungen wurden neben der Aufgabenstellung auch eine Referenzlösung, ein
Startprogramm und Beispielwörter benötigt.
Die KI nutzte ich vor allem, um Varianten bestehender Aufgaben zu erzeugen,
beispielsweise mit anderem Alphabet, anderen Wortlängen oder unterschiedlichen
Schwierigkeitsgraden.

Ob eine Aufgabe tatsächlich korrekt war, habe ich anschließend mit der automatischen
Übungskontrolle des Simulators gegen die Referenzlösung geprüft
(siehe Abschnitt~\ref{4_2_EA}).
Nur geprüfte Aufgaben wurden in den Aufgabenkatalog aufgenommen
(siehe Abbildung~\ref{fig:ki_aufgaben}).

\begin{figure}[htbp]
	\centering
	\begin{tikzpicture}[
		font=\footnotesize,
		box/.style={draw, rounded corners=2pt, align=center, inner sep=5pt, minimum height=1.4cm},
		flow/.style={-{Stealth[length=2.2mm,width=1.8mm]}, semithick},
		lbl/.style={font=\scriptsize\itshape, inner sep=2pt}
	]
		\node[box, text width=2.6cm, fill=black!5]  (a) at (0,0)    {Bestehende Aufgabe\\als Vorlage};
		\node[box, text width=3.4cm, fill=black!5]  (b) at (4.1,0)  {KI erzeugt Variante\\(Alphabet, Wortlänge, Schwierigkeit)};
		\node[box, text width=3.2cm, fill=black!12] (c) at (8.6,0)  {Übungskontrolle prüft gegen die Referenzlösung};
		\node[box, text width=2.6cm, fill=black!5]  (d) at (13.2,0) {Aufnahme in den Aufgabenkatalog};

		\draw[flow] (a) -- (b);
		\draw[flow] (b) -- (c);
		\draw[flow] (c) -- (d) node[lbl, midway, above=2pt, fill=white] {bestanden};
		\draw[flow] (c.south) -- ++(0,-1.0) coordinate (fb) -| (b.south);
		\node[lbl, below=1mm] at ($(fb)!0.5!(fb -| b.south)$)
			{Abweichung: Variante wird verworfen};
	\end{tikzpicture}
	\caption[Erzeugung und Prüfung von Aufgabenvarianten]{%
		Erzeugung und Prüfung von Aufgabenvarianten. Nur Varianten, die sich
		gegenüber der Referenzlösung identisch verhalten, werden in den
		Aufgabenkatalog übernommen.}
	\label{fig:ki_aufgaben}
\end{figure}

\subNsection{2}{Kritische Reflexion}{Kristian Popov}

Die KI lieferte häufig brauchbare, aber nicht immer korrekte Vorschläge.
Teilweise verwendete sie APIs, die in der verwendeten Version nicht mehr existierten,
oder erfand Funktionen und Bibliotheken.
Auch Sonderfälle und Auswirkungen auf andere Teile des Projekts wurden gelegentlich
übersehen.

Ein weiteres Problem war die starke Ähnlichkeit zwischen den Themen
\emph{Endliche Automaten} und \emph{Reguläre Sprachen}.
Statt bestehende Funktionen zu abstrahieren, erzeugte die KI teilweise doppelte
Implementierungen.
Änderungen, die mehrere Themen gleichzeitig betrafen, musste ich deshalb häufig in
kleinere Schritte aufteilen.

Insgesamt sparte mir die KI besonders bei wiederkehrenden Aufgaben und beim Ausbau
bestehender Funktionen viel Zeit.
Bei theoretischen Konzepten und größeren Architekturänderungen war jedoch weiterhin
meine eigene Prüfung entscheidend.
Ein Orchestrator wie Sakana Fugu könnte künftig die Auswahl des passenden Modells
übernehmen~\cite{sakana2026fugu}; da das System erst kurz vor Abschluss meiner Arbeit
erschien, konnte ich es nicht mehr einsetzen.

Für mich war die KI damit vor allem ein Werkzeug zur Unterstützung.
Die fachlichen Entscheidungen und die Verantwortung für die Korrektheit der
Implementierung lagen jederzeit bei mir.



	\subNsection{0}{Designentscheidungen der Fachmodule}
\subNsection{1}{Turing Maschine}{Author}

Die Turing-Maschine ist als eigenständiger Themenbereich innerhalb des Moduls
\glqq Grundlagen der Informatik\grqq{} umgesetzt. Der Bereich besteht nicht nur
aus einem Simulator, sondern aus einem vollständigen Lern-Workflow mit
Übersicht, Übungen, Beispielen und einer freien Simulation. Die Git-Historie
zeigt dabei eine schrittweise Entwicklung: Zuerst wurden Datenmodell,
Simulator-Schnittstelle und Debugger aufgebaut, danach folgten die editierbare
Konfiguration, die responsive Banddarstellung, Beispiele und Übungen,
Internationalisierung, Import/Export, automatische Speicherung und zuletzt
weitere Tests sowie die Integration in gemeinsam genutzte Komponenten.

\subNsection{2}{Entscheidungen in der Architektur}{Author}

\begin{itemize}
  \item Die Implementierung trennt die Turing-Maschine konsequent in
  \texttt{domain}, \texttt{application}, \texttt{presentation} und
  \texttt{adapters}. Dadurch liegen fachliche Regeln, React-Zustand,
  UI-Komponenten und die Abbildung auf generische Tabellenkomponenten in
  getrennten Schichten.

  \item Das statische Maschinenprogramm und der Laufzeitzustand sind getrennte
  Modelle. \texttt{TMProgram} beschreibt Zustände, Alphabete, Blank-Symbol,
  Startzustand, Haltezustände, Übergänge und Anfangsband. \texttt{TMRuntime}
  beschreibt aktuellen Zustand, Kopfposition, Bandinhalt, Schrittzähler und
  Halteinformation.

  \item Die eigentliche Simulation liegt in reinen Funktionen wie
  \texttt{createInitialRuntime}, \texttt{readCurrentSymbol} und
  \texttt{stepRuntime}. Diese Funktionen hängen nicht von React, i18n oder der
  Darstellung ab und sind deshalb gut isoliert testbar.

  \item Fehlende Übergänge werden fachlich als normales Halten modelliert und
  nicht als technischer Fehler. Ein fehlender Übergang erzeugt einen
  \texttt{haltReason} vom Typ \texttt{missing\_transition}; ein regulärer
  Haltezustand erzeugt einen \texttt{haltReason} vom Typ
  \texttt{halt\_state}. Dadurch kann die UI sauber zwischen einem erwartbaren
  Verhalten einer partiell definierten Turing-Maschine und echten
  Implementierungsfehlern unterscheiden.

  \item Die Validierung prüft Struktur und Konsistenz, aber nicht die
  Vollständigkeit der Übergangsfunktion. Es wird unter anderem geprüft, ob
  Start- und Haltezustände existieren, Symbole nur aus einzelnen Zeichen
  bestehen, das Blank-Symbol gültig ist und Übergänge auf existierende
  Zustände und Bandsymbole verweisen. Unvollständige Tabellen bleiben dagegen
  erlaubt, weil sie didaktisch und fachlich sinnvoll sind.

  \item Für die Simulation wird eine Klasse \texttt{TuringMachineSimulator}
  verwendet, die den gemeinsamen \texttt{ISimulator}-Vertrag der Anwendung
  erfüllt. Intern nutzt sie aber weiterhin die reinen Domain-Funktionen. Damit
  bleibt die Turing-Maschine kompatibel zur allgemeinen Simulationsengine,
  ohne die Fachlogik in Hooks oder UI-Komponenten zu verschieben.

  \item Der Simulator speichert eine Historie der Laufzeitzustände. Dadurch
  kann die UI neben \glqq Schritt vorwärts\grqq{} auch \glqq Schritt zurück\grqq{}
  anbieten. Die Historie liegt in der Simulator-Anwendungsschicht und bleibt
  von der Darstellung unabhängig.

  \item Die Editorlogik ist in \texttt{tmEditorReducer} und
  \texttt{useTMConfigEditor} ausgelagert. Der Reducer kümmert sich um
  Operationen wie Zustände hinzufügen, Zustände umbenennen, Bandsymbole
  umbenennen, Übergänge aktualisieren und Haltezustände verwalten. Wichtig ist
  dabei, dass Referenzen mitgeführt werden: Wird ein Zustand umbenannt, werden
  auch Startzustand, Haltezustände und Zielzustände in Übergängen angepasst.

  \item Die Übergangstabelle ist intern als verschachtelte Struktur aus
  Lesesymbol und Zustand modelliert. Für die Darstellung wird sie über
  \texttt{tmTableAdapter} in Header, Zeilen und aktive Tabellenzelle
  übersetzt. Diese Adapterentscheidung verhindert, dass die Domain-Schicht
  Tabellen- oder UI-Wissen enthält.

  \item Eingabealphabet und Bandalphabet werden nicht vollständig manuell
  gepflegt. Das Eingabealphabet wird aus dem Anfangsband abgeleitet; das
  Bandalphabet wird aus vorhandenem Bandalphabet, Anfangsband, Blank-Symbol,
  Lesesymbolen und Schreibsymbolen abgeleitet. Dadurch bleibt die
  Konfiguration für Lernende einfacher, während die Domain dennoch beide
  Konzepte kennt.

  \item Der Import und Export ist als eigener Transfer-Codec umgesetzt. Die
  aktuelle Architektur verwendet eine Bitstream-basierte, Base64-URL-sichere
  Kodierung mit Symbolregistern, Roh-UTF-8-Fallback und verschiedenen
  Übergangstabellen-Kodierungen. Die Historie zeigt, dass dies aus einem
  einfachen Base64-Proof-of-Concept zu einer kompakten, robusteren
  Speicherform weiterentwickelt wurde.

  \item Für gespeicherte Aufgabenstände wird der gemeinsame
  Training-Center-Speicher genutzt. Beispiele werden bewusst nicht dauerhaft
  überschrieben, während freie Simulationen und Übungen eigene
  Speicherkennungen verwenden. Dadurch bleiben eingebaute Beispiele
  reproduzierbar und Übungsfortschritte lokal erhalten.

  \item Die Übungsprüfung ist von der sichtbaren Simulation getrennt. Für jeden
  Validierungsfall wird ein Simulator mit dem jeweiligen Eingabewort gestartet,
  bis zum Halten oder bis zu einer Schrittgrenze ausgeführt und dann mit der
  erwarteten Ausgabe verglichen. Zusätzlich reicht eine passende Ausgabe nicht:
  Die Maschine muss in einem gültigen Haltezustand enden.
\end{itemize}

\subNsection{2}{Entscheidungen im UI}{Author}

\begin{itemize}
  \item Die Turing-Maschine nutzt den gemeinsamen Themenaufbau der Plattform:
  Kopfbereich mit Titel und Beschreibung, Tabs für Übersicht, Übungen,
  Beispiele und Simulation sowie die globale Sidebar mit Sprach- und
  Theme-Umschaltung.

  \item Die Simulation startet im Konfigurationsmodus. Lernende sehen zuerst
  Anfangsbandwort, Startzustand, Übergangstabelle, Haltezustände,
  Eingabealphabet, Bandalphabet sowie Import/Export. Erst nach dem Anwenden
  der Konfiguration wechselt die Ansicht in den Laufzeitmodus.

  \item Die Übergangstabelle ist der zentrale Editor. Zeilen stehen für
  Bandsymbole, Spalten für Nicht-Haltezustände. Haltezustände werden bewusst
  nicht als editierbare Übergangsspalten angezeigt, weil von ihnen aus keine
  weiteren Übergänge erwartet werden.

  \item Übergänge werden in der Form
  \texttt{(Next State, Write Symbol, Movement)} eingegeben. Das UI erlaubt
  leere Zellen und behandelt sie als nicht definierte Übergänge. Damit passt
  die Darstellung zur Architekturentscheidung, partielle Übergangsfunktionen
  zuzulassen.

  \item Fehler in einzelnen Tabellenzellen blockieren den Editor gezielt. Bei
  ungültigen Symbolen oder ungültigem Übergangsformat bleibt die betroffene
  Zelle editierbar, andere Änderungen und der Start der Simulation werden
  deaktiviert. So wird verhindert, dass ein halbvalider Tabellenzustand
  unbemerkt übernommen wird.

  \item Das Blank-Symbol wird sichtbar als \(\diamond\) dargestellt. Für die
  Eingabe ist zusätzlich die Tastaturform \texttt{<>} vorgesehen, die in das
  Blank-Symbol umgewandelt wird. Das reduziert Eingabehürden, ohne in der
  Darstellung ein uneindeutiges Leerzeichen verwenden zu müssen.

  \item Die Laufzeitansicht besteht aus Band, Steuerleiste, Statuskarte und
  Übergangstabelle. Das Band zeigt ein Fenster um die Kopfposition mit
  Indexbeschriftungen und einer Pfeilmarkierung für den aktuellen Kopf. Die
  Tabellenzelle des aktuellen Zustands und Symbols kann über den Adapter als
  aktive Zelle markiert werden.

  \item Das Band ist responsiv ausgelegt. Die Anzahl sichtbarer Bandzellen wird
  aus der verfügbaren Breite berechnet; Animationen werden nur für kleine
  Verschiebungen und sinnvolle Geschwindigkeiten ausgeführt. Dadurch bleibt
  die Darstellung auch bei kleinen Viewports und schnellen Simulationen
  beherrschbar.

  \item Die Steuerung verwendet die gemeinsamen Simulator-Controls der
  Plattform: Play, Pause, Schritt, Schritt zurück, Reset und
  Geschwindigkeitsregler. Damit verhält sich die Turing-Maschine ähnlich wie
  andere Simulatoren der Anwendung.

  \item Die Statuskarte zeigt aktuellen Zustand, aktuelles Symbol,
  Schrittzahl, Haltestatus und bei fehlenden Übergängen eine fachliche
  Erklärung. Bei der Live-Prüfung hielt die Maschine ohne definierten Übergang
  mit der Meldung, dass für den aktuellen Zustand und das aktuelle Symbol kein
  Übergang definiert ist.

  \item Übungen erweitern die Simulation um die Aktionen
  \glqq Lösung laden\grqq{} und \glqq Lösung prüfen\grqq{}. Das Prüfergebnis
  wird in einem Dialog als Tabelle mit Eingabe, Ausgabe, erwarteter Ausgabe,
  Bestehen und gültigem Haltezustand angezeigt.

  \item Die Oberfläche ist weitgehend lokalisiert. Eine sichtbare Auffälligkeit
  in der deutschen UI ist jedoch, dass in der Übergangstabelle der
  Zeilen-Hinzufügen-Text noch als \glqq Add new row\grqq{} erscheint. Das
  deutet auf einen fehlenden oder nicht angebundenen Übersetzungseintrag in
  einer gemeinsam genutzten Tabellenkomponente hin.
\end{itemize}

\subNsection{2}{Andere Entscheidungen}{Author}

\begin{itemize}
  \item Die Beispiele und Übungen sind als Szenarien modelliert. Beispiele
  enthalten ein fertiges Programm; Übungen enthalten zusätzlich ein
  Starterprogramm, ein Lösungsprogramm, eine Speicher-ID und mehrere
  Validierungsfälle. Damit kann derselbe Simulator für freie Arbeit, Beispiele
  und bewertbare Übungen genutzt werden.

  \item Die enthaltenen Beispiele und Übungen decken verschiedene typische
  Aufgaben ab, etwa Bit-Invertierung, unäres und binäres Inkrementieren,
  Löschen des Bandes, Suchen von Mustern, Duplizieren von Nullen,
  Binäraddition und einfache Worterkennung. Die Übungen sind nach
  Schwierigkeitsgrad sortiert.

  \item Die Testsuite ist für diesen Bereich vergleichsweise breit. Es gibt
  Unit-Tests für Domain-Logik, Formatierung, Programmvalidierung,
  Transfer-Codec, Editor-Reducer, Tabellenadapter, Bandlogik,
  Lösungskontrolle und Simulator-Hook. Die E2E-Tests prüfen zusätzlich die
  sichtbare Bedienung, Beispiel- und Übungsnavigation, Local-Storage-Verhalten,
  Import/Export, Fehlervalidierung und das Erstellen einer Bit-Inverter-Maschine
  von Grund auf.

  \item Die Git-Historie legt nahe, dass die Architektur bewusst nachträglich
  stabilisiert wurde. Nach einem ersten Datenmodell und Simulator folgten
  größere Refactorings nach einer Architekturplanung, danach responsive
  Verbesserungen, Internationalisierung, Test-IDs, Übungsintegration und
  persistente Speicherung. Das spricht für eine iterative Entwicklung, bei der
  die UI zunächst funktionsfähig gemacht und anschließend stärker in Schichten
  getrennt wurde.

  \item Besonders sinnvoll ist die Entscheidung, \glqq fehlender Übergang\grqq{}
  nicht als Programmfehler zu behandeln. Für Lernende ist eine unvollständige
  Turing-Maschine ein normaler Zwischenzustand beim Konstruieren. Gleichzeitig
  kann die Übungsprüfung streng bleiben, weil sie nur dann besteht, wenn die
  Maschine nicht nur die richtige Ausgabe produziert, sondern auch in einem
  vorgesehenen Haltezustand endet.

  \item Ein technischer Zielkonflikt liegt im Transferformat. Die kompakte
  Bitstream-Kodierung ist gut für kurze Exportstrings und lokale Speicherung,
  erschwert aber die manuelle Lesbarkeit. Diese Entscheidung ist plausibel,
  weil der Import/Export im UI als Kopier- und Austauschfunktion gedacht ist
  und nicht als primäres Bearbeitungsformat.

  \item Ein weiterer Zielkonflikt liegt in der abgeleiteten Alphabetpflege.
  Sie reduziert Eingabeaufwand und verhindert viele Inkonsistenzen, kann aber
  für fortgeschrittene theoretische Diskussionen weniger explizit sein als ein
  separat editierbares Eingabealphabet. Die aktuelle Lösung wirkt für die
  Lernplattform angemessen, weil sie den Editor einfacher hält.

  \item Aus Wartungssicht ist positiv, dass die meisten UI-Elemente stabile
  \texttt{data-testid}-Attribute besitzen. Das erleichtert E2E-Tests und war
  auch bei der Live-Inspektion hilfreich. Gleichzeitig sollten kleine
  Lokalisierungslücken in gemeinsam genutzten Komponenten weiter beobachtet
  werden, damit deutsche und englische Oberfläche konsistent bleiben.
\end{itemize}


\subNsection{1}{Endliche Automaten}{Kristian Popov} \label{4_2_EA}

Im Thema \emph{Endliche Automaten} können Studierende deterministische und
nichtdeterministische endliche Automaten selbst definieren und Eingabewörter Schritt
für Schritt ausführen.
Jeder Simulationsschritt wird gleichzeitig in der Übergangstabelle und im
Zustandsdiagramm dargestellt.
Ergänzt wird die Einzelsimulation durch einen Stapeltest für mehrere Wörter, die
Umwandlung rechtslinearer Grammatiken in Automaten und Übungen mit automatischer
Rückmeldung.

\subNsection{2}{Entscheidungen in der Architektur}{Kristian Popov}

Ursprünglich war nur die Unterstützung deterministischer Automaten vorgesehen.
Nichtdeterministische kamen hinzu, weil dort ein Symbol in mehrere Zustände zugleich
führen darf: Ein Automat ist so oft direkt aus der Sprachbeschreibung ablesbar und
kompakter, und die Mehrdeutigkeit muss erst beim Schritt zum DFA aufgelöst werden.
Da jeder DFA gleichzeitig ein NFA ist, wurde für beide Varianten bewusst nur eine
Ausführungslogik verwendet.
Der Kern arbeitet mit einer Menge aktiver Zustände, die bei einem DFA höchstens ein
Element enthält; $\varepsilon$-Übergänge werden dabei stets über die
$\varepsilon$-Hülle aufgelöst.
Der deklarierte Automatentyp steuert Validierung und Eingabemaske, nicht aber die
Ausführung.

Der Simulator ist in fünf Bereiche aufgeteilt:

\begin{description}[leftmargin=*,style=nextline,itemsep=0.75ex]
\item[\texttt{domain}]
enthält die Fachlogik wie Simulation, Validierung, Grammatikumwandlung, Stapeltest
und Übungskontrolle.

\item[\texttt{application}]
verbindet die Fachlogik mit React und verwaltet Simulation und Editor.

\item[\texttt{adapters}]
bereitet die Daten für die Übergangstabelle auf.

\item[\texttt{presentation}]
enthält die Benutzeroberfläche mit Diagramm, Tabelle, Eingabewort, Stapeltest und
Grammatikeingabe.

\item[\texttt{config}]
enthält Standardautomaten und Beispiele.
\end{description}

Die Trennung sorgt dafür, dass die Fachlogik unabhängig von React bleibt und ohne
gerenderte Oberfläche testbar ist.
Das Thema \emph{Reguläre Sprachen} folgt demselben Schnitt und greift zusätzlich
direkt auf die Grammatikumwandlung dieses Themas zurück.

Ein Simulationsschritt erzeugt einen neuen Zustand mit der bisherigen Historie.
Dadurch sind sowohl das Zurückgehen als auch die Auswahl früherer Schritte möglich,
ohne eine zusätzliche Historie in der Oberfläche zu verwalten.

Auch die Validierung ist von der Oberfläche getrennt.
Sie liefert Fehlercodes und relevante Werte, während die Oberfläche daraus die
passenden Meldungen erzeugt.
Dadurch bleibt die Fachlogik unabhängig von der Sprache der Benutzeroberfläche.

Die mitgelieferten Automaten sind bewusst partiell definiert.
Fehlt ein Übergang, wird das Wort abgelehnt.
Auf einen zusätzlichen Fangzustand wurde verzichtet, da dieser zwar die Sprache nicht
verändert, aber die Diagramme unnötig vergrößert.

\subNsection{2}{Entscheidungen im UI}{Kristian Popov}

Ursprünglich war für die Simulation nur die Übergangstabelle vorgesehen.
Während der Entwicklung zeigte sich jedoch, dass der Ablauf eines Wortes damit schwer
nachzuvollziehen ist: Zyklen und Sackgassen werden erst sichtbar, wenn man mehrere
Zeilen gedanklich verknüpft.
Deshalb wurde zusätzlich ein Zustandsdiagramm integriert.

Beide Darstellungen sind an denselben Simulationsschritt gebunden: Die Tabelle hebt
genau die Zelle hervor, die dem im Diagramm markierten Übergang entspricht.
Sie zeigt die formale Definition des Automaten, während das Diagramm den
tatsächlichen Weg des Wortes sichtbar macht.
Bei größeren Automaten wird das automatisch berechnete Layout schnell
unübersichtlich.
Einzelne Zustände lassen sich deshalb frei ziehen, der Ausschnitt vergrößern und
verkleinern und die Ausgangsansicht mit einem Klick wiederherstellen.
Aktive und fehlerhafte Übergänge bleiben dabei auch bei geringer Zoomstufe
beschriftet.

Verstöße gegen den Determinismus werden nicht nur aufgelistet, sondern unmittelbar an
Zustand und Übergang im Diagramm markiert.
Damit ist die Ursache dort sichtbar, wo sie entsteht.

Der Editor umfasst alle wesentlichen Bestandteile eines Automaten: Automatentyp,
Alphabet, Zustände, Start- und Endzustände sowie das Eingabewort.
Die Zustandsmenge wird direkt über die Übergangstabelle verwaltet.
Alphabetsymbole sind auf einzelne Zeichen beschränkt, damit eine Eingabe wie
\verb|10| nicht versehentlich als ein einzelnes Symbol aufgenommen wird.
Die Zeile für $\varepsilon$-Übergänge erscheint nur bei einem NFA, sobald sie belegt
ist oder ausdrücklich angefordert wird.

\subNsection{2}{Andere Entscheidungen}{Kristian Popov}

Neben der Einzelsimulation gibt es einen Schrittverlauf, über den bereits erreichte
Schritte direkt ausgewählt werden können.
Der gewählte Schritt wird dabei vom Anfang des Laufs an neu berechnet statt
zwischengespeichert; bei den hier üblichen Wortlängen ist das unkritisch und hält
Anzeige und tatsächlichen Ablauf zwingend deckungsgleich.
Der Stapeltest ermöglicht es, mehrere Wörter gleichzeitig zu prüfen.
Ein Ergebnis kann anschließend direkt in die Einzelsimulation übernommen werden.

Eine weitere Funktion ist die Umwandlung rechtslinearer Grammatiken in Automaten.
Nichtterminale werden dabei zu Zuständen und Produktionen werden in entsprechende
Übergänge oder Endzustände übersetzt.
Für die Eingabe werden zwei Schreibweisen unterstützt, die automatisch erkannt
werden.

Beispiele und Übungen sind getrennt aufgebaut.
Beispiele enthalten fertige Automaten und dienen zum Nachvollziehen, während Übungen
mit einem leeren Automaten beginnen und zusätzlich Aufgabenstellung, Referenzlösung
und automatische Kontrolle enthalten.
Die Übungen sind nach \emph{Sprache} $\rightarrow$ \emph{Automat} und
\emph{Grammatik} $\rightarrow$ \emph{Automat} gegliedert.

Die Übungskontrolle vergleicht nicht die Struktur der Automaten, sondern deren
Verhalten.
Dazu werden alle Wörter über dem gemeinsamen Alphabet bis zu einer festgelegten Länge
geprüft.
Der Automat der Studierenden muss für jedes Wort dieselbe Entscheidung treffen wie
die Referenzlösung.
Bei einer Abweichung wird das erste unterschiedliche Wort als Gegenbeispiel
angezeigt.
Dadurch können auch unterschiedlich aufgebaute, aber funktional korrekte Automaten
akzeptiert werden.
Die Prüfung stellt jedoch keinen formalen Äquivalenznachweis dar.
Strukturelle Mängel wie ungültige Übergänge werden vorrangig gemeldet, da sie meist
die eigentliche Ursache sind und eine Liste fehlgeschlagener Testwörter dann in die
Irre führen würde.

Der Bearbeitungsstand jeder Übung wird gespeichert und beim erneuten Öffnen
wiederhergestellt, sodass der Wechsel zwischen Übungen keine bereits begonnene
Bearbeitung überschreibt.


\subNsection{1}{Reguläre Sprachen}{Kristian Popov}

Das Thema \emph{Reguläre Sprachen} unterstützt die Arbeit mit regulären Sprachen
und rechtslinearen Grammatiken.
Es bietet drei miteinander verknüpfte Funktionen:
Die Aufzählung der von einer Grammatik erzeugten Wörter, die Erzeugung einer
passenden Grammatik aus einer Sprachbeschreibung sowie die Prüfung einzelner Wörter.
Sprachbeschreibungen können entweder als Aufzählung oder in Mengenschreibweise
eingegeben werden.
Bei der Aufzählung einer Sprache werden die Wörter nach ihrer Länge sortiert.
Die Ausgabe ist dabei auf eine maximale Wortlänge und Wortanzahl begrenzt.
Bei der Wortprüfung zeigt das Thema neben der Annahme oder Ablehnung auch den Grund
für ein abweichendes Ergebnis an, beispielsweise einen nicht möglichen Übergang.
Beide Umwandlungsrichtungen sind zusätzlich als Übungen mit automatischer Rückmeldung verfügbar.

\subNsection{2}{Entscheidungen in der Architektur}{Kristian Popov}

Der Simulator besteht aus einer koordinierenden Komponente und fachlich getrennten Bausteinen.
Die Komponente \texttt{RLSimulator} verwaltet die drei Arbeitsmodi,
den Wechsel zwischen ihnen und den Einstieg in Übungen.
Sie enthält selbst keine formale Logik.
Diese liegt vollständig im Verzeichnis \texttt{domain} und ist damit
unabhängig von React testbar.

Die fachlichen Bausteine sind wie folgt aufgeteilt:

\begin{description}[leftmargin=*,style=nextline,itemsep=0.75ex]
	\item[Grammatikverarbeitung]
		Der Grammatikparser und die Validierung lesen die Eingabe ein, prüfen die
		Struktur einer rechtslinearen Grammatik und überführen sie in gemeinsame
		Datentypen.

	\item[Sprachbeschreibung]
		Ein eigener Parser verarbeitet Aufzählungen und Mengenschreibweisen.
		Er liefert eine strukturierte Beschreibung aus Faktoren, Exponenten und
		Nebenbedingungen, die anschließend in eine Grammatik umgewandelt werden
		kann.

	\item[Umwandlung und Vereinfachung]
		Die Umwandlung erzeugt aus einer Sprachbeschreibung eine strikt
		rechtslineare Grammatik.
		Ergänzende Funktionen fassen gleichwertige Nichtterminale zusammen und
		formatieren das Ergebnis für die Ausgabe.

	\item[Sprachauswertung]
		Die Aufzählung erzeugt Wörter bis zu den festgelegten Grenzen; die
		Worterkennung prüft ein einzelnes Wort gegen eine Grammatik.
		Die Übungsprüfung baut auf denselben Funktionen auf und ergänzt sie um die
		Auswertung der vorgegebenen Testfälle.
\end{description}

Wiederverwendet werden der Grammatikparser, die Grammatikvalidierung und die gemeinsamen
Datentypen.
Eine eigene Ableitungslogik direkt auf den Produktionen hätte zwei
Implementierungen erfordert, die dasselbe Verhalten abbilden müssen.
Mit der gemeinsamen Grundlage stimmen Aufzählung, Worterkennung und Übungsprüfung
stets überein.

Für die Sprachbeschreibung habe ich einen eigenen Parser geschrieben, der zwei
Formen versteht:

\begin{description}[leftmargin=*,style=nextline,itemsep=0.75ex]
	\item[Aufzählung]
		listet die Wörter unmittelbar auf:\\
		\verb|L = {"010", "0010"}|

	\item[Mengenschreibweise]
		beschreibt ein Muster mit Variablen und Nebenbedingungen:\\
		\verb+L = {a^n c b^m | n, m > 0}+
\end{description}

Endet eine Aufzählung mit einer Fortsetzungsmarkierung, versucht der Parser, daraus
ein wiederkehrendes Muster abzuleiten, statt die Wörter nur wörtlich zu übernehmen.
So lässt sich eine unendliche Sprache angeben, ohne die Mengenschreibweise schon
beherrschen zu müssen.

Bestimmte Eingaben werden bewusst abgelehnt, auch wenn sie sich formal beschreiben lassen.
Verwendet dieselbe Variable mehrere Faktoren, wie in \verb|a^n b^n|, koppelt das
deren Längen aneinander.
Der Musterfall \verb|a^n b^n| ist nicht regulär; die unterstützte Sprachbeschreibung
lässt solche gekoppelten Längenbedingungen deshalb generell nicht zu und weist sie
mit einer eigenen Meldung zurück.
Obere Schranken werden ebenfalls nicht unterstützt,
obwohl sie endliche und damit reguläre Sprachen beschreiben.
Sie würden die Eingabesprache unnötig erweitern und lassen sich klarer als Aufzählung angeben.
Diese Fälle ausdrücklich zu erklären, ist lehrreicher, als sie stillschweigend zu verwerfen.

Die erzeugte Grammatik ist strikt rechtslinear, es entstehen also ausschließlich
Regeln der Form \verb|A -> aB| und \verb|A -> a| sowie die Regel
$A \rightarrow \varepsilon$.
Eine frühere Fassung gab eine kompaktere Form wie \verb|S -> A c B| aus.
Diese Schreibweise ist zwar kürzer, aber eine allgemeine kontextfreie Regel und
verlässt damit die Typ-3-Grammatiken, um die es in diesem Thema gerade geht.

Um trotzdem eine lesbare Ausgabe zu bieten, gibt es zwei Sichten auf dasselbe
Ergebnis statt zweier Formalismen:
Die detaillierte Sicht zeigt ein Nichtterminal je Präfix, die vereinfachte Sicht
führt Nichtterminale zusammen, die dieselben Wörter ableiten.
Bei der vereinfachten Sicht werden äquivalente Teilbäume zusammengeführt.
Sie bleibt rechtslinear und kann die Grammatik nur verkleinern, nie ihre Sprache verändern.

Alle Stellen, die eine rechtslineare Grammatik ausgeben, benutzen dieselbe
Namensvergabe und Ausgabefunktion.
Zuvor lag diese Logik in drei leicht unterschiedlichen Kopien vor.
In einer davon verbrauchte ein weggelassener Endzustand einen Buchstaben, sodass in der Ausgabe
eine Lücke entstand und etwa \verb|E| übersprungen wurde.
Eine gemeinsame Funktion schließt solche Abweichungen aus und stellt sicher, dass
eine exportierte Grammatik nach dem Import dieselben Namen trägt.

\subNsection{2}{Entscheidungen im UI}{Kristian Popov}

Die drei Arbeitsweisen sind über eine Modusauswahl erreichbar.
Ihre Komponenten bleiben im Hintergrund erhalten und werden nur ein- oder ausgeblendet.
Dadurch bleiben Eingaben beim Wechsel zwischen den Modi bestehen.
Würden die Modi jeweils neu aufgebaut, ginge die begonnene Arbeit verloren
und die Ansicht spränge auf ihr Standardbeispiel zurück.

Ist eine Übung geöffnet, ersetzt eine geführte Arbeitsfläche den freien Umwandler.
Vorgabe und Aufgabe stehen dabei nebeneinander:
In der einen Richtung ist die Grammatik gegeben und die Sprache gesucht, in der
anderen umgekehrt.
Ein Moduswechsel während einer offenen Übung würde den aktuellen Bearbeitungsstand
verwerfen. Deshalb wird vorher eine Bestätigung verlangt.

Fehlermeldungen erscheinen unmittelbar unter dem jeweiligen Eingabefeld, aber erst,
nachdem überhaupt etwas eingegeben wurde.
Ein leeres Feld ist kein Fehler, sondern der Ausgangszustand.
Zusätzlich lässt sich die unterstützte Syntax direkt neben dem Eingabefeld
aufklappen, damit die Schreibweise nicht erst gesucht werden muss.

\subNsection{2}{Andere Entscheidungen}{Kristian Popov}

Die Kontrolle der Übungen vergleicht keine Texte.
Für die Richtung \emph{Sprache}$\rightarrow$ \emph{Grammatik} sind zu jeder Aufgabe Wörter hinterlegt,
die erzeugt werden müssen, und Wörter, die nicht erzeugt werden dürfen.
Geprüft wird, ob die eingegebene Grammatik alle geforderten Wörter ableitet und
alle ausgeschlossenen ablehnt.
Ein Textvergleich mit einer Musterlösung wäre hier unbrauchbar, da dieselbe Sprache
von vielen unterschiedlichen Grammatiken erzeugt werden kann und die Sprachen zum
Teil unendlich sind.
Die Rückmeldung nennt daher falsch erzeugte und fehlende Wörter statt nur einer
Ja-Nein-Aussage.
Da nur vorbereitete Beispiele geprüft werden, dient auch diese Kontrolle
der gezielten Rückmeldung und ist kein allgemeiner Beweis für die Gleichheit zweier Sprachen.

Für die Richtung \emph{Grammatik} $\rightarrow$ \emph{Sprache} verwendet die
Kontrolle dieselben Grenzen für Wortlänge und Wortanzahl wie die Schaltfläche zum
Anzeigen der Lösung.
Andernfalls könnte eine angezeigte Lösung die anschließende Prüfung nicht bestehen.

Beim Import und Export unterscheiden sich die beiden Formate bewusst.
Eine Grammatik wird strukturiert kodiert, also in Nichtterminale, Terminale und
Produktionen zerlegt und bitweise gepackt.
Eine Sprachbeschreibung wird dagegen als Text übernommen, weil sie vom Benutzer
selbst formuliert ist und ihre Schreibweise erhalten bleiben soll.
Beide Seiten prüfen beim Kodieren und beim Dekodieren, dass der Inhalt gültig ist.

Jeder Modus speichert seinen Stand in einem eigenen Bereich der lokalen Ablage.
Ohne diese Trennung hätten sich die Grammatik des Aufzählungsmodus und die des
Worterkenners gegenseitig überschrieben, da beide auf demselben Eingabetyp arbeiten.

Die Aufgabenstellungen selbst sind in einer sprachneutralen Notation hinterlegt und
nicht in den Übersetzungsdateien.
Eine Produktion wie \verb|START = "1", A| ist in beiden Oberflächensprachen
identisch; eine Übersetzung würde hier nur zusätzliche Fehlerquellen schaffen.


\subNsection{1}{Zahlensysteme}{Enrico Ebert}

\subNsection{2}{Entscheidungen in der Architektur}{Enrico Ebert}

\subNsection{2}{Entscheidungen im UI}{Enrico Ebert}

\subNsection{2}{Andere Entscheidungen}{Enrico Ebert}

Die Punkte sind nur Vorschläge ...

\subNsection{1}{Moore- und Mealy-Automaten}{Markus Nöth}

Das Modul vermittelt den Entwurf synchroner Schaltwerke mit Tabellen, Zustandsgraphen und Logikschaltungen. Beide Automatentypen bieten freie Simulationen, Beispiele und Übungen.

\subNsection{2}{Entscheidungen in der Architektur}{Markus Nöth}

Die Implementierung beider Automatentypen baut auf dem gemeinsamen Interface \texttt{ISimulator\allowbreak<TConfig,\allowbreak{}TState>} der Plattform auf. Die mathematische Definition und die Aufgabengestaltung orientieren sich an den Vorlesungsfolien von Prof.~Braun sowie am Lehrbuch von Vossen und Witt~\cite[Kapitel~4]{vossen2000}: Beim Moore-Automaten hängt die Ausgabe nur vom aktuellen Zustand ab ($\lambda: S \to \Gamma$). Beim Mealy-Automaten ist die Ausgabefunktion an das Wertepaar aus aktuellem Zustand und anliegendem Eingabesymbol gekoppelt ($\lambda: S \times \Sigma \to \Gamma$).

Die Logik liegt in den Klassen \verb|MooreMachineSimulator| und \verb|MealyMachineSimulator|. Bei jedem Taktschritt sichern sie den gesamten Ausführungszustand in einer Historie ab. Dadurch kann der Benutzer über \verb|stepBack()| deterministisch zurückspringen, indem der gesicherte Zustand wiederhergestellt wird.

Für beide Automatentypen erzeugt das Modul Wahrheitstabellen, disjunktive Normalformen (\gls{dnf}) und Logikgatter:
\begin{itemize}
  \item \textbf{Moore-Transformation (\texttt{buildMooreConversion}):} Da die Ausgabe allein an den Zustand gekoppelt ist, hängen die Ausgangsvariablen ($y$) nur von den Zustandsvariablen ($S_0 \dots S_n$) ab. Die Übergangsfunktionen für den Folgezustand ($T_0 \dots T_n$) berechnen sich aus den Zustands- und Eingangsvariablen ($E_0 \dots E_m$).
  \item \textbf{Mealy-Transformation (\texttt{buildMealyConversion}):} Sowohl der Folgezustand als auch die Ausgänge hängen gleichzeitig von Zustand und Eingabe ab. Die Wahrheitstabelle der Ausgänge umfasst daher die Bitbelegungen der definierten Zustands-Eingabe-Paare.
\end{itemize}
Aus den generierten Wahrheitstabellen berechnet das angebundene DNF-Modul über den Funktionsaufruf \texttt{compute\-Smallest\-Boolean\-Expression()} die minimierten Schaltfunktionen. Bis vier Variablen wird das globale Literalminimum für UND/ODER-Formeln mit negierten Literalen bestimmt. Bei fünf bis acht Variablen wird zwischen minimaler DNF und ihrer Faktorisierung gewählt, ohne globale Optimalitätsgarantie. Darüber bleibt die kanonische DNF. Die Grenze begrenzt den Aufwand, garantiert aber keine blockierungsfreie Berechnung.


\begin{figure}[htbp]
  \centering
  \begin{tikzpicture}[
    font=\footnotesize,
    box/.style={draw, rounded corners=2pt, align=center, inner sep=5pt, text width=4.2cm, minimum height=0.8cm},
    flow/.style={-{Stealth[length=2mm]}, semithick}
  ]
    \node[box, fill=black!5] (config) at (0,0) {Editor-Konfiguration};
    \node[box, fill=black!12] (sim) at (-3.5,-1.5) {Übernommene Konfiguration\\Simulator mit Historie};
    \node[box, fill=black!12] (conv) at (3.5,-1.5) {Konvertierung\\Wahrheitstabellen und Formeln};
    \node[box] (graph) at (-3.5,-3) {Laufzeitansicht\\Tabelle und Zustandsgraph};
    \node[box] (logic) at (3.5,-3) {Konvertierungsansicht\\Tabellen und Logikgatter};
    \draw[flow] (config.south) -- ++(0,-0.25) -| (sim.north);
    \draw[flow] (config.south) -- ++(0,-0.25) -| (conv.north);
    \draw[flow] (sim) -- (graph);
    \draw[flow] (conv) -- (logic);
  \end{tikzpicture}
  \caption{Die Simulation übernimmt die geprüfte Konfiguration; die Konvertierung nutzt den Editorstand.}
  \label{fig:moore_mealy_datenfluss}
\end{figure}

\subNsection{2}{Entscheidungen im UI}{Markus Nöth}

Die Konfiguration der Automaten erfolgt über Editoren, deren Erscheinungsbild für beide Typen bewusst einheitlich gehalten ist:
\begin{itemize}
  \item \textbf{Editor-Tabellen:} Bei Mealy-Automaten werden Folgezustand und Ausgabesymbol gemeinsam in einer Matrixzelle editiert. Bei Moore-Automaten sind Übergangsmatrix und Zustandsausgabetabelle in zwei getrennte Bereiche aufgeteilt.
  \item \textbf{Transformationsansichten:} In der Konvertierungsansicht (\verb|AutomataConversionView|) kann der Anwender zwischen Übergangs- und Ausgabetabellen wechseln. Jede Zeile schlüsselt die symbolischen Werte in ihre Bitbelegung auf. Die zugehörige Formelliste zeigt kanonische und minimierte Ausdrücke für jedes Bit an.
  \item \textbf{Zustandsgraph (\texttt{StateMachineGraph}):} Diese Komponente visualisiert den Automaten als gerichteten Graphen. Sie verteilt Knoten, bündelt gleichgerichtete Kanten und versetzt Gegenrichtungen. Selbstschleifen erscheinen als Kurven. Zoom und Pan verändern den sichtbaren Koordinatenausschnitt über die \gls{svg}-\verb|viewBox|. Hervorgehoben werden der aktuelle Zustand und die Kante für das nächste Eingabesymbol.
  \item \textbf{Logikgatter-Diagramm (\texttt{LogicGateDiagram}):} Stellt wahlweise die kanonischen oder minimierten Schaltfunktionen für jedes Bit mit Gattersymbolen für AND, OR und NOT dar. Knoten und einzelne Verbindungen lassen sich durch Anklicken hervorheben.
\end{itemize}
Beide Komponenten wurden modular unter \verb|app/src/components/shared/| abgelegt, sodass sie neben Moore- und Mealy-Maschinen auch von künftigen Simulatoren der Plattform direkt wiederverwendet werden können.

\subNsection{2}{Andere Entscheidungen}{Markus Nöth}

Das Modul enthält elf praxisnahe Übungsaufgaben: sieben für Moore-Automaten sowie vier für Mealy-Automaten. Jede Aufgabe stellt ein Starterprogramm, eine Lösung und Validierungsfälle bereit.

Für Speicherung und Austausch kommt der Modul-Codec \verb|programTransfer.ts| zum Einsatz: Anstelle von unkomprimiertem \gls{json} serialisieren ein \verb|BitWriter| und \verb|BitReader| Zustände, Symbole, Übergänge und Ausgaben direkt in dichte Bitfolgen. Dies spart den Syntax-Overhead von Attributnamen ein: Bezeichner werden nur einmalig deklariert und in den Übergängen als Bit-Indizes minimaler Breite ($\lceil\log_2 N\rceil$) referenziert, während Kanten über 1-Bit-Flags kodiert werden. Die Bitfolgen werden Base64URL-kodiert und ermöglichen kompakte \glspl{url} für das Teilen von Automaten sowie die persistente Ablage im Local Storage des Trainingscenters.



\subNsection{1}{Disjunktive Normalform}{Duy Tan Pham}

Die \gls{dnf} ist ein eigener Themenbereich im Modul
\glqq Grundlagen der Informatik\grqq{}. Sie folgt dem gleichen Aufbau wie die anderen
Themen: Übersicht, Übungen, Beispiele und eine freie Simulation. Die Inhalte bauen
schrittweise aufeinander auf: von der kanonischen DNF über die minimale DNF bis zur
optional global kleinsten Formel. Die einzelnen Schritte der Ableitung werden dabei
sichtbar gemacht.

Die inhaltliche Ausrichtung folgt den vorhandenen Übungsunterlagen der Lehrveranstaltung.
Dort wird die DNF aus Wahrheitstabellen mit drei und vier Variablen gebildet. Der Modus
\glqq Wahrheitstabelle nach DNF\grqq{} setzt genau diesen Aufgabentyp interaktiv um. Die
Grenze von vier Variablen deckt damit zugleich die vorgesehenen Übungsbeispiele ab.

\subNsection{2}{Entscheidungen in der Architektur}{Duy Tan Pham}

Abbildung~\ref{fig:dnf_ablauf} fasst den fachlichen Datenfluss der beiden
Eingabemodi zusammen.

\begin{figure}[htbp]
  \centering
  \begin{tikzpicture}[
    font=\footnotesize,
    box/.style={draw, rounded corners=2pt, align=center, inner sep=5pt, minimum height=1.0cm},
    flow/.style={-{Stealth[length=2.2mm,width=1.8mm]}, semithick},
    bound/.style={flow, dashed},
    lbl/.style={font=\scriptsize\itshape, inner sep=2pt, fill=white}
  ]
    \node[box, text width=2.1cm, fill=black!5]  (source)    at (0,0)
      {Boolescher\\Ausdruck};
    \node[box, text width=2.1cm, fill=black!5]  (parser)    at (3,0)
      {Parser und\\Syntaxbaum};
    \node[box, text width=2.1cm, fill=black!12] (table)     at (6,0)
      {Referenz-\\tabelle};
    \node[box, text width=2.1cm, fill=black!5]  (canonical) at (9,0)
      {Kanonische DNF};
    \node[box, text width=2.1cm, fill=black!12] (minimal)   at (12,0)
      {Minimale DNF\\exakte Minimierung};

    \node[box, text width=2.1cm, fill=black!5]  (tableinput) at (0,-2.3)
      {Vorgegebene\\Tabelle};
    \node[box, text width=2.1cm, fill=black!5]  (answer)     at (3,-4.6)
      {Antwort der\\Lernenden};
    \node[box, text width=2.5cm, fill=black!12] (check)      at (6,-4.6)
      {Prüfung von Form\\und Äquivalenz};
    \node[box, text width=2.7cm, fill=black!12] (global)     at (10.5,-3.0)
      {Optional: global\\kleinste Formel};

    \draw[flow] (source) -- (parser);
    \draw[flow] (parser) -- (table);
    \draw[flow] (table) -- (canonical);
    \draw[flow] (canonical) -- (minimal);
    \draw[flow] (tableinput.east) -- (1.5,-2.3) -| (table.south)
      node[lbl, near start, above] {direkt};
    \draw[flow] (answer) -- (check);
    \draw[flow] (table.south) -- (check.north)
      node[lbl, midway, right] {Referenz};
    \draw[flow] (table.south east) to[out=-25,in=165]
      node[lbl, midway, above] {Funktion} (global.west);
    \draw[bound] (minimal.south) to[out=-90,in=35]
      node[lbl, midway, right] {Obergrenze} (global.north east);
  \end{tikzpicture}
  \caption[Fachlicher Datenfluss im DNF-Bereich]{%
    Fachlicher Datenfluss im DNF-Bereich. Die Referenztabelle verbindet beide
    Eingabemodi mit Ableitung, Minimierung, globaler Suche und Äquivalenzprüfung.}
  \label{fig:dnf_ablauf}
\end{figure}

\begin{itemize}
  \item Der Bereich trennt \texttt{domain}, \texttt{application}, \texttt{adapters}
  und \texttt{presentation}. Parser, Wahrheitstabelle, DNF-Ableitung, Minimierung und
  Prüfung liegen in \texttt{domain} und lassen sich ohne React testen.

  \item Der boolesche Ausdruck wird von einem eigenen Parser (\texttt{dnfParser.ts}) in
  einen Syntaxbaum umgewandelt. Aus diesem Baum wird die Wahrheitstabelle erzeugt. Die
  Wahrheitstabelle ist die gemeinsame Grundlage für alle weiteren Schritte. Der Parser
  erlaubt bewusst mehrere Schreibweisen (\texttt{!}, \texttt{\&}, \texttt{|} und ebenso
  \(\lnot\), \(\land\), \(\lor\)) und meldet Fehler mit Position.

  \item Die Ableitung trennt drei Ergebnisstufen sauber: die kanonische DNF
  (\texttt{dnfDerivation.ts}, ein Vollminterm je Eins-Zeile \cite[S.~163]{schiffmann2003}), die minimale DNF
  (\texttt{quineMcCluskey.ts}) und die optional global kleinste Formel
  (\texttt{minimalExpression.ts}). Berechnung und Darstellung bleiben dabei getrennt.

  \item Die wichtigste fachliche Entscheidung betrifft die Minimierung. Die erste Version
  wählte die restliche Überdeckung nach Quine-McCluskey
  \cite[S.~172]{schiffmann2003} \emph{greedy} aus und konnte dadurch die beste Lösung
  verpassen. Quine-McCluskey bildet nun die Primimplikanten. Die Funktion
  \texttt{selectCover} wählt daraus mit Branch-and-Bound exakt eine vollständige
  Abdeckung mit minimaler Literalzahl. Zweige, die das beste bekannte Ergebnis nicht
  mehr verbessern können, werden frühzeitig verworfen.

  \item Für die optional kleinste Formel wird zusätzlich eine exakte Synthese über den
  Funktionsraum berechnet. Die minimale DNF liefert dafür eine gültige Obergrenze. Für
  Ausdrücke aus Literalen, die mit UND und ODER verbunden werden, findet die Suche bis
  vier Variablen das globale Literalminimum. Negationen stehen dabei nur direkt an
  Variablen. Da die Anzahl möglicher Funktionen mit \(2^{2^{n}}\) stark wächst, ist das
  Werkzeug auf vier Variablen begrenzt. Diese Grenze wird in Anwendungsschicht,
  Tabellen-Eingabe und Import geprüft.

  \item Die Prüfung ist von der Darstellung getrennt. Ein eingegebener Ausdruck wird für
  jede Zeile der Wahrheitstabelle ausgewertet und nicht durch einen reinen Textvergleich
  bewertet. Dadurch erkennt die Übungsprüfung (\texttt{dnfExerciseCheck.ts}) korrekte,
  nur äquivalente und nicht äquivalente Antworten. Die Vereinfachungsprüfung
  (\texttt{dnfSimplificationCheck.ts}) prüft zusätzlich, ob die Literalzahl das exakte
  Minimum erreicht.

  \item Import, Export und automatisches Speichern verwenden den gemeinsamen
  Training-Center-Speicher und einen kompakten, URL-sicheren Base64-Code.
\end{itemize}

\subNsection{2}{Entscheidungen im UI}{Duy Tan Pham}

\begin{itemize}
  \item Der DNF-Bereich bietet zwei Aufgabenmodi: \glqq Ausdruck nach
  Wahrheitstabelle\grqq{} und den umgekehrten Weg \glqq Wahrheitstabelle nach DNF\grqq{}.
  Beim zweiten Modus stellen die Lernenden aus einer gegebenen Tabelle die DNF selbst auf.
  Beide Übungen nutzen dieselben Bausteine.

  \item Die einzelnen Schritte der Ableitung werden sichtbar gemacht: von den Eins-Zeilen
  der Tabelle über die passenden kanonischen Terme bis zur minimalen Form. Bei der
  Vereinfachung wird nicht nur das Ergebnis gezeigt, sondern auch der Lösungsweg. In den
  Übungen bearbeiten die Lernenden Wahrheitstabellen oder DNF-Ausdrücke selbst und
  erhalten eine fachlich berechnete Rückmeldung.

  \item Kleinere Wahrheitstabellen werden vollständig angezeigt. Bei \texttt{16} Zeilen,
  also bei vier Variablen, gibt es eine kompakte, scrollbare Ansicht, die sich auf eine
  Vollansicht umschalten lässt. Dadurch bleibt die maximale Darstellungshöhe begrenzt.

  \item Folgeaktionen bleiben sichtbar und werden bis zum Erfüllen ihrer Voraussetzungen
  ausgegraut. Nur kontextabhängige Schaltflächen, etwa für die Vollansicht einer großen
  Tabelle, erscheinen bei Bedarf. Dadurch bleibt die Bedienung vorhersehbar.

  \item Eingaben werden vereinheitlicht: Lernende dürfen \texttt{!}, \texttt{\&} und
  \texttt{|} tippen, angezeigt wird aber immer \(\lnot\), \(\land\), \(\lor\). Das macht die
  Eingabe leichter, ohne die Darstellung uneinheitlich zu machen.

  \item Der DNF-Bereich nutzt den gemeinsamen Themenaufbau mit Kopfbereich, Tabs und
  globaler Sidebar. Dadurch fügt er sich in die übrige Lernplattform ein.
\end{itemize}

\subNsection{2}{Andere Entscheidungen}{Duy Tan Pham}

Die kanonische und die minimale DNF werden bewusst getrennt dargestellt. In der
kanonischen DNF enthält jeder Minterm alle Variablen. Eine minimale DNF kann dagegen
kürzere Terme besitzen, etwa \((a \land c) \lor (b \land c)\). Die getrennte Darstellung
macht diesen fachlichen Unterschied nachvollziehbar und verhindert, dass eine korrekte
Vereinfachung fälschlich als unvollständig angesehen wird.


\subNsection{1}{Encoding}{Enrico Ebert}

Ziel des Encoding-Moduls ist es, die Transformation zwischen unstrukturiertem Text, Codepoints und binären Byte-Sequenzen transparent zu machen. Wir haben die praxisrelevanten Standards UTF-8 und Base64 vollständig implementiert. Der Fokus liegt darauf, die mathematischen und bitweisen Zwischenschritte – wie Bitmasken bei UTF-8 oder das Padding bei Base64 – für Studierende visuell nachvollziehbar darzustellen.

\subNsection{2}{Architektur}{Enrico Ebert}

Die Architektur trennt Datenmodell, Transformationslogik und UI strikt. Der Zustand wird zentral verwaltet. Child-Komponenten besitzen keinen eigenen Zustand für verarbeitete Daten, sondern geben Änderungen über State Lifting an die Simulator-Komponente weiter. Um beim Wechsel von Übungen einen sauberen Zustand ohne Artefakte zu garantieren, erzwingen wir über dynamische \texttt{key}-Attribute gezielte React-Remounts.

Die UTF-8-Logik nutzt konsequent \texttt{Array.from()} zur sicheren Iteration über Unicode-Codepoints, um Fehler durch Surrogate-Pairs (UTF-16) zu vermeiden. Der Decoder validiert Byteströme streng und weist Sicherheitsrisiken wie \textit{Overlong Encodings} konsequent ab. Die Base64-Pipeline arbeitet iterativ in vier Stufen: ASCII-Konvertierung, 8-Bit-Verkettung, 6-Bit-Aufteilung und abschließendes Padding (\texttt{======}).

Übungen werden programmgesteuert dynamisch generiert; eine statische Liste existiert nicht. Der Generator erzeugt pro Schwierigkeitsgrad (z.\,B. reines ASCII für \textit{easy}, 3-Byte-Unicode für \textit{hard}) Aufgaben inklusive der kanonischen Lösung. Da diese Aufgaben ephemere IDs erhalten, wird der Fortschritt hier im Gegensatz zum freien Modus nicht dauerhaft gespeichert.

\subNsection{2}{UI-Entscheidungen}{Enrico Ebert}

Das UI visualisiert abstrakte Bitoperationen detailliert auf Byte-Ebene. Jeder Zwischenschritt lässt sich aufklappen. Bei UTF-8 werden Steuerbits zur besseren Unterscheidung rot und Nutzbits blau hervorgehoben. Bei Base64 weisen wir die 6-Bit-Gruppen und das Padding explizit aus. Nicht-druckbare Zeichen werden als Escape-Sequenzen (z.\,B. \texttt{\textbackslash x0A}) dargestellt, um das Layout nicht zu brechen.

Wir verzichten bewusst auf zeitgesteuerte Animationen. Der statische \textit{Step-by-Step}-Ansatz gibt den Studierenden die nötige Zeit, die Transformationen in ihrem eigenen Tempo zu analysieren.

Im Übungsmodus gibt ein Validierungspanel die Lösung erst bei korrekter Eingabe frei. Ein Warn-Dialog verhindert versehentlichen Fortschrittsverlust beim Tab-Wechsel. Das automatische Speichern ist im Übungsmodus deaktiviert und greift nur im freien Simulator-Modus.

\subNsection{2}{Andere Entscheidungen}{Enrico Ebert}

Der Modul-übergreifende Zustandstransfer (\textit{Import/Export}) erfolgt Unicode-sicher via Base64URL-Kodierung inklusive Fallback auf Standard-JSON.


\subNsection{1}{Java-Programmierumgebung}{Markus Nöth}

Die Java-Programmierumgebung ist Bestandteil des Moduls \glqq Programmieren~1\grqq{} und ermöglicht das Schreiben, Kompilieren und Ausführen von Java-Programmen direkt im Browser ohne Server-Backend.

\subNsection{2}{Entscheidungen in der Architektur}{Markus Nöth}

Die Umgebung folgt einer mehrschichtigen Architektur aus Präsentation, Orchestrierung und isolierter Ausführungsumgebung:
\begin{enumerate}
  \item \textbf{Präsentationsschicht:} Der \verb|JavaExecutionProvider| stellt den vom Hook verwalteten Zustand als \verb|JavaExecutionSnapshot| bereit (Status, Konsolenausgaben, Diagnosen und Artefakte). Im Monaco-Editor~\cite{monacoeditor} (\verb|JavaEditorPanel|) begrenzt die Anwendung den Quelltext über \verb|javaEditorLimits.ts| auf 500.000~UTF-16-Codeeinheiten und 10.000~Zeilen und visualisiert Fehlermeldungen direkt im Code.
  \item \textbf{Orchestrierungsschicht:} Der Hook \verb|useJavaExecution| steuert den Lebenszyklus des Web Workers (\verb|javaRunner.worker.ts|) über asynchrone Nachrichten mit eindeutigen Request-\glspl{id}. Eine konfigurierbare Laufzeitüberwachung schützt vor Endlosschleifen: Überschreitet der Auftrag einschließlich Laden, Kompilieren und Ausführen das Zeitlimit, beendet der Hauptthread den Worker mittels \texttt{worker.\allowbreak{}terminate()} und startet eine neue Instanz. Die Programmausführung blockiert damit nicht den Hauptthread.
  \item \textbf{Compiler- und \gls{wasm}-Laufzeit:} Ein untergeordneter Compiler-Worker lädt das vorkompilierte \gls{wasm}-Compilermodul (\texttt{teavm-javac}) samt Klassenbibliotheken als statische Assets nach. Dieses bündelt das OpenJDK-\verb|javac|-Frontend mit dem TeaVM-Backend~\cite{teavm}: Nach der Syntax- und Typprüfung durch \verb|javac| übersetzt TeaVM~\cite{teavm} den erzeugten Bytecode in \gls{wasm}, das anschließend im Runner-Worker ausgeführt wird.
\end{enumerate}



Abbildung~\ref{fig:java_worker_ablauf} zeigt den Ablauf:
\begin{enumerate}
  \item \textbf{Quelltext-Vorverarbeitung:} Echte interaktive Konsoleneingaben während der Programmausführung würden eine zusätzliche Abstimmung zwischen Benutzeroberfläche und Wasm-Laufzeit erfordern. Aus diesem Grund setzt das Modul auf einen leichtgewichtigen Vorprozessor: \verb|inlineScannerInput()| ersetzt erkannte Scanner-Aufrufe vor der Kompilierung per regulärem Ausdruck durch feste Werte aus einem vorbereiteten \verb|ScannerInputCursor|. Unterstützt werden unter anderem ganze Zahlen (\texttt{nextInt()}), Dezimalzahlen (\texttt{nextDouble()}), Wahrheitswerte (\texttt{nextBoolean()}) sowie einzelne Wörter und Textzeilen (\texttt{next()}, \texttt{nextLine()}). Bei gültigen, vorab bereitgestellten Eingaben lassen sich damit einfache Programme zur Eingabe, Berechnung und Ausgabe umsetzen, die die grundlegenden Arbeitsschritte für den Einstieg abdecken. Die Ersetzung folgt jedoch der Reihenfolge im Quelltext: Wiederholte Eingaben in Schleifen oder nur bedingt ausgeführte Eingabeaufrufe werden nicht entsprechend dem tatsächlichen Kontrollfluss ausgewertet.
  \item \textbf{Übersetzung und Ausführung:} Der vorverarbeitete Quelltext wird im Web Worker durch \verb|javac| kompiliert und via TeaVM~\cite{teavm} in ein \gls{wasm}-Modul transformiert, das im Runner-Worker instanziiert wird.
  \item \textbf{Pufferung des Ausgabestroms:} Über das Foreign Function Interface (\gls{ffi}) von WebAssembly registriert \verb|installTeaVmConsole()| die Host-Funktionen \texttt{putcharStdout} und \texttt{putcharStderr} in der Import-Tabelle des Wasm-Moduls. Standard- und Fehlerausgaben von \texttt{System.out} bzw. \texttt{System.err} werden so zeichenweise über die FFI-Grenze an JavaScript übergeben. \texttt{create-Buffered-Worker-Output} bündelt die Zeichen je Ausgabestrom mit einer konfigurierten Verzögerung von 50ms. Der Puffer wird gesendet, sobald der Timer ausgeführt werden kann, bei einem weiteren Zeichen nach Ablauf der Verzögerung oder beim Abschluss der Ausführung. Dies reduziert die Zahl einzelner \texttt{postMessage}-Nachrichten und damit auch die Zahl der Rerenders.
\end{enumerate}


\begin{figure}[H]
  \centering
  \begin{tikzpicture}[
    font=\footnotesize,
    box/.style={draw, rounded corners=2pt, align=center, inner sep=5pt, text width=3.8cm, minimum height=1cm},
    flow/.style={-{Stealth[length=2mm]}, semithick},
    lbl/.style={font=\scriptsize, fill=white, inner sep=2pt}
  ]
    \node[box, fill=black!5] (ui) at (0,0) {Hauptthread\\Editor und Ausführungs-Hook};
    \node[box, fill=black!12] (runner) at (5,0) {Runner-Worker\\Scanner-Vorverarbeitung};
    \node[box, fill=black!5] (compiler) at (10,0) {Compiler-Worker\\javac $\rightarrow$ Bytecode\\TeaVM $\rightarrow$ Wasm};
    \node[box, fill=black!12] (wasm) at (5,-2) {Runner-Worker\\Wasm-Ausführung und Ausgabepuffer};
    \draw[flow] (ui) -- (runner);
    \draw[flow] (runner) -- (compiler);
    \draw[flow] (compiler.south) |- (wasm.east) node[lbl, near end, above] {Wasm};
    \draw[flow] (wasm.west) -| (ui.south) node[lbl, near start, above] {Ausgaben, Status};
  \end{tikzpicture}
  \caption[Java-Ausführung mit zwei Workern]{Java-Ausführung mit getrenntem Runner- und Compiler-Worker und TeaVM~\cite{teavm} als Wasm-Backend.}
  \label{fig:java_worker_ablauf}
\end{figure}

\subNsection{2}{Entscheidungen im UI}{Markus Nöth}

Die Arbeitsoberfläche orientiert sich an vertrauten \glspl{ide} wie Visual Studio Code:
\begin{itemize}
  \item \textbf{Monaco-Editor:} Der Editor bietet Code-Faltung, Volltextsuche (\texttt{Strg+F}) und automatische Einrückung. Compilerfehler mit Positionsangabe, beispielsweise fehlende Semikolons oder inkompatible Typen, erscheinen als rote Marker. Entfernte Scanner-Zeilen werden dabei nicht auf den ursprünglichen Quelltext zurückgerechnet.
  \item \textbf{Eingabebereich für \texttt{System.in}:} Über ein separates Texteingabefeld kann der Benutzer Konsoleneingaben vorab definieren, die der Vorprozessor beim Programmstart verarbeitet.
  \item \textbf{Konsole und Problemliste (\texttt{JavaConsolePanel}):} Standardausgabe, Fehlerausgabe und Diagnosen erscheinen in getrennten Abschnitten. Die Diagnoseliste nennt die Meldung sowie, falls vorhanden, Datei und Position; ein Sprung in den Editor ist dort nicht implementiert.
  \item \textbf{Steuerleiste:} Start und Abbruch liegen beim Simulator; Zurücksetzen und Fokusmodus bietet der gemeinsame Simulationsrahmen.
\end{itemize}

\subNsection{2}{Andere Entscheidungen}{Markus Nöth}

Das Modul enthält vier interaktive Beispielprogramme und 13~Übungsaufgaben (u.\,a. Konsolenausgaben, Schleifen, Variablenberechnungen).
Die Aufgabenprüfung kompiliert den eingegebenen Code und vergleicht seine normalisierte Standardausgabe mit den erwarteten Ausgaben definierter Testfälle. Dies belegt keine allgemeine Korrektheit. Bei Aufgaben mit vorgegebenen Variablenbelegungen ersetzt ein regulärer Ausdruck (\verb|sourceTransform|) die Initialisierungswerte im Code durch alternative Testdaten, um weitere Eingabefälle zu prüfen. Ändert der Lernende jedoch die Bezeichner oder den vom Muster erwarteten Aufbau dieser Vorgabezeilen, greift der Ersetzungsausdruck nicht mehr. In diesem Fall kann die Validierung fehlschlagen, obwohl der geschriebene Algorithmus fachlich korrekt ist. Eine Lösung für dieses Problem wäre das gezielte Sperren von Zeilen im Editor (siehe Abschnitt~\ref{6_2_Perspektiven}).

Für Speicherung und Austausch wählt \verb|javaCodeTransfer.ts| den kürzesten Transfer-String aus drei Verfahren: direkte Wörterbuch-Tokens, deflatierte Tokens oder deflatiertes UTF-8, jeweils gekennzeichnet durch ein Formatpräfix (\texttt{A}, \texttt{B} oder \texttt{C}) mit Base64URL-Nutzdaten. Das integrierte Wörterbuch ersetzt häufige Java-Idiome (z.\ B. \texttt{System.out.println} oder Methodensignaturen) sowie Umlaute durch kompakte 1-Byte-Tokens im Bereich oberhalb des Standard-ASCII-Codes. Falls die Browser-Laufzeit keine Deflate-Kompression unterstützt, stehen die unkomprimierten Tokens als Fallback bereit. Der Export validiert den Quelltext strikt und akzeptiert nur die im Codec zugelassenen Zeichen.



	\subNsection{0}{Evaluation \& Projekterfahrung}
\subNsection{1}{Qualitätssicherung}{Enrico Ebert} \label{5_1_QS}

CI/CD-Pipeline (GitLab Runner auf Raspberry Pi), automatisiertes Linting, Formatierung (Prettier), Unit-Tests (Vitest) und E2E-Tests (Playwright)


\subNsection{1}{Abgleich mit dem Lastenheft}{Kristian Popov}

Wurden die Qualitätsziele (Performance, Usability, Barrierefreiheit durch Dark Mode und i18n) erreicht?

\subNsection{1}{Design-Iterationen}{Duy Tan Pham}

Welche Entscheidungen mussten im Prozess revidiert werden?


	\subNsection{0}{Fazit \& Ausblick}
\subNsection{1}{Zusammenfassung der Projektergebnisse}{Alle}

Kristian Popov: Mein Schwerpunkt im Projekt lag auf den Themen
\emph{Endliche Automaten} und \emph{Reguläre Sprachen}. Für Endliche Automaten
habe ich eine interaktive Umgebung umgesetzt, in der deterministische und
nichtdeterministische Automaten erstellt und Eingabewörter schrittweise ausgeführt
werden können. Übergangstabelle und Zustandsdiagramm zeigen dabei denselben Ablauf,
sodass die formale Definition und ihre grafische Darstellung miteinander verbunden
werden. Ein Stapeltest prüft mehrere Wörter auf einmal und macht dadurch sichtbar, ob
ein Automat die beabsichtigte Sprache erkennt und nicht nur das gerade ausprobierte
Wort. Ergänzt wird das Thema durch Beispiele, Übungen mit Rückmeldung und die
Umwandlung rechtslinearer Grammatiken in Automaten.

Im Thema Reguläre Sprachen stehen drei Arbeitsweisen zur Verfügung: Studierende können
die von einer Grammatik erzeugten Wörter aufzählen, aus einer Sprachbeschreibung eine
passende Grammatik erzeugen und einzelne Wörter auf ihre Zugehörigkeit zur Sprache
prüfen. Bei der Wortprüfung wird neben dem Ergebnis auch dessen Grund angezeigt.
Übungen ermöglichen auch hier eine direkte Überprüfung der eigenen Lösung. In beiden
Themen lassen sich Automaten, Grammatiken und Sprachbeschreibungen über die gemeinsame
Import- und Exportfunktion der Plattform weitergeben. Bei beiden Themen war mir wichtig,
die Fachlogik klar von der Benutzeroberfläche zu trennen. Dadurch konnten gemeinsame
Grammatik- und Automatenfunktionen wiederverwendet werden, während die einzelnen
Funktionen testbar und übersichtlich blieben.

Duy Tan Pham: Mein fachlicher Schwerpunkt lag auf dem Themenbereich
\emph{Disjunktive Normalform}. Dort können Studierende boolesche Ausdrücke und
Wahrheitstabellen untersuchen, die kanonische DNF schrittweise herleiten und Ausdrücke
vereinfachen. Beispiele zeigen einen vollständigen Lösungsweg. In den Übungen bearbeiten
die Lernenden dagegen Wahrheitstabellen oder DNF-Ausdrücke selbst und erhalten eine
fachlich berechnete Rückmeldung. Dabei werden auch logisch gleichwertige Schreibweisen
erkannt.

Die Minimierung liefert eine DNF mit minimaler Literalzahl. Zusätzlich kann für bis zu
vier Variablen eine global kleinste Formel aus UND- und ODER-Verknüpfungen bestimmt
werden. Die Fachlogik ist von der Benutzeroberfläche getrennt und durch automatisierte
Tests abgesichert. Der DNF-Bereich nutzt außerdem den gemeinsamen Themenaufbau sowie die
vorhandenen Funktionen für Import und Export. Dadurch fügt er sich in Bedienung und
Darstellung in die gesamte Lernplattform ein.


\subNsection{1}{Perspektiven}{Kristian Popov} \label{6_2_Perspektiven}

Die modulare Architektur erlaubt eine flexible Erweiterung der Anwendung. Künftige Perspektiven gliedern sich in folgende Schwerpunkte:

\begin{itemize}
    \item \textbf{Weitere Lernbereiche und Topics:} Neue Themen – etwa von anderen Professoren beigesteuerte Fachmodule – lassen sich über den zentralen Modulkatalog dynamisch in Navigation und Routing integrieren.
    \item \textbf{Aufgabenverwaltung und -varianten:} Lehrende können Aufgaben und Vorlagen direkt über die Oberfläche anlegen, automatisch Varianten erzeugen sowie Quizfragen und Lernkarten ergänzen.
    \item \textbf{Mobile Nutzung:} Optimierung der Benutzeroberfläche – insbesondere von Übergangstabellen und Zustandsdiagrammen – für kleine Bildschirme und Tablets.
    \item \textbf{Java-Zeilensperrung:} Abgrenzung bearbeitbarer Code-Bereiche im Editor durch Marker, um nicht zu bearbeitende Zeilen schreibgeschützt zu halten.
    \item \textbf{Gamification und LMS-Anbindung:} Motivationselemente wie Streaks oder Abzeichen sowie eine Moodle-Anbindung zur zentralen Lernstandsauswertung.
    \item \textbf{Offene Punkte aus dem Lastenheft (Abschnitt~\ref{5_2_LH}):}
    \begin{itemize}
        \item \textit{Usability-Erhebung:} Durchführung strukturierter Tests zur Bedienbarkeit über einen internen Testzugang.
        \item \textit{Barrierefreiheit \& Mobil-Profiling:} Automatisierte CI-Prüfungen für Kontraste, Tastaturbedienung und mobile Performance-Messungen.
    \end{itemize}
\end{itemize}

Während sich neue Lernbereiche, Aufgabenvarianten und Editor-Sperren vollständig in die bestehende, wartungsfreie Client-Architektur integrieren lassen, würden zentrale Auswertungen, Nutzerkonten oder eine direkte Moodle-Anbindung ein Backend erfordern.
\end{Content}

\begin{PostContent}
	\subNsection*[toc]{0}{Anhang}

	\subNsection*[toc]{1}{Abkürzungsverzeichnis}

	% Nimmt alle in glossary.tex definierten Akronyme auf, auch wenn sie im Text
	% nicht per \gls{...} ausgezeichnet sind.
	\glsaddall

	% Ohne diesen Schalter setzt der list-Stil zwischen den Anfangsbuchstaben
	% jeweils einen Absatzabstand, wodurch die Liste auseinandergezogen wirkt.
	\glsnogroupskiptrue

	\printacronyms[heading=none]

	\subNsection*[toc]{1}{Literatur}

	\printbibliography[heading=none]
\end{PostContent}

\end{document}
